%!TEX root = ../main.tex

\chapter*{Nomenclatura}
\addcontentsline{toc}{chapter}{Nomenclatura}
\label{ch:nomenclatura}

En este capítulo se define el convenio de ejes y los ángulos de actitud sobre los
que se construye el PFD, y se recogen las siglas y acrónimos utilizados a lo largo del proyecto. 

\section{Convenio de ejes y actitud}
\label{sec:nom_ejes}

A lo largo del proyecto se emplea el convenio aeronáutico de ejes NED (\textit{North-East-Down}),
en el que los ejes de navegación apuntan al Norte, al Este y hacia abajo
(en el sentido de la gravedad). El sistema de ejes solidario al cuerpo
(\textit{body frame}) tiene su origen en la IMU, con el eje $x$ hacia el morro
(adelante), el eje $y$ hacia el costado derecho y el eje $z$ hacia abajo.

La actitud estimada por el filtro de fusión se expresa mediante tres ángulos de
Euler en el convenio NED, que se corresponden con los tres grados de libertad de
orientación de la aeronave:

\begin{itemize}
    \item \textbf{Cabeceo} (\textit{pitch}): rotación alrededor del eje $y$
    (morro arriba o abajo).
    \item \textbf{Alabeo} (\textit{roll}): rotación alrededor del eje $x$
    (inclinación lateral de las alas).
    \item \textbf{Rumbo} (\textit{heading}): rotación alrededor del eje $z$
    (orientación respecto al Norte magnético).
\end{itemize}

En las figuras y en el código estos tres ángulos se abrevian como P, R y H,
respectivamente. El cuaternión de actitud $\mathbf{q}$ es la representación
interna de la orientación, a partir de la cual se obtienen los ángulos de Euler.

\newpage
\section{Siglas y acrónimos}
\label{sec:nom_acronimos}

Las siglas y acrónimos propios del ámbito aviónico, de sensores o específicos del proyecto se
recogen en la \cref{tab:nom_acronimos}.

\begin{xltabular}{\linewidth}{@{} l X @{}}
    \caption{Siglas y acrónimos utilizados en el proyecto.} \label{tab:nom_acronimos} \\
    \toprule
    \textbf{Sigla} & \textbf{Significado} \\
    \midrule
    \endfirsthead
    % Cabecera de las siguientes páginas
    \toprule
    \textbf{Sigla} & \textbf{Significado} \\
    \midrule
    \endhead
    % Pie de tabla
    \bottomrule
    \endlastfoot
    % Contenido
    AHRS    & \textit{Attitude and Heading Reference System} (sistema de referencia de actitud y rumbo) \\
    ARINC   & \textit{Aeronautical Radio, Incorporated} \\
    ASA     & \textit{Sensitivity Adjustment} (valores de ajuste de sensibilidad de fábrica del magnetómetro AK8963) \\
    BNR     & \textit{Binary Number Representation} (representación numérica binaria en ARINC 429) \\
    CEP     & \textit{Circular Error Probable} (error circular probable) \\
    EKF     & \textit{Extended Kalman Filter} (filtro de Kalman extendido) \\
    EMA     & \textit{Exponential Moving Average} (media móvil exponencial) \\
    FMCW    & \textit{Frequency-Modulated Continuous Wave} (onda continua modulada en frecuencia) \\
    GNSS    & \textit{Global Navigation Satellite System} (sistema global de navegación por satélite) \\
    ICGC    & \textit{Institut Cartogràfic i Geològic de Catalunya} (proveedor de teselas de mapa) \\
    IIR     & \textit{Infinite Impulse Response} (respuesta al impulso infinita) \\
    IMU     & \textit{Inertial Measurement Unit} (unidad de medida inercial) \\
    ISM     & \textit{Industrial, Scientific and Medical} (banda de radio de uso libre) \\
    LDO     & \textit{Low-Dropout} (regulador lineal de baja caída de tensión) \\
    MARG    & \textit{Magnetic, Angular Rate and Gravity} (sensor de campo magnético, velocidad angular y gravedad) \\
    MSL     & \textit{Mean Sea Level} (nivel medio del mar, referencia de altitud) \\
    NED     & \textit{North-East-Down} (convenio de ejes Norte-Este-Abajo) \\
    NMEA    & \textit{National Marine Electronics Association} (formato de tramas de los receptores GNSS) \\
    PFD     & \textit{Primary Flight Display} (pantalla primaria de vuelo) \\
    SDI     & \textit{Source/Destination Identifier} (identificador de fuente y destino en ARINC 429) \\
    SSM     & \textit{Sign/Status Matrix} (matriz de signo y estado en ARINC 429) \\
    ToF     & \textit{Time of Flight} (tiempo de vuelo) \\
    USB-CDC & \textit{USB Communications Device Class} (puerto serie virtual sobre USB) \\
    VSI     & \textit{Vertical Speed Indicator} (indicador de velocidad vertical) \\
\end{xltabular}

\null